\section{Public Key Infrastructure (PKI)}
Masalah utama kriptografi asimetris adalah: \textit{Bagaimana kita tahu bahwa kunci publik yang kita gunakan benar-benar milik orang yang bersangkutan?}

\subsection{Otoritas Sertifikat (Certificate Authority - CA)}
CA adalah pihak ketiga terpercaya yang menerbitkan sertifikat digital. Sertifikat ini mengikat identitas seseorang (misal: nama perusahaan) dengan kunci publiknya.

\subsection{Sertifikat X.509}
X.509 adalah standar format sertifikat digital yang mencakup:
\begin{itemize}
    \item Versi sertifikat.
    \item Nomor seri.
    \item Nama penerbit (CA).
    \item Masa berlaku (valid from/to).
    \item Nama pemilik sertifikat (Subject).
    \item Kunci publik pemilik.
    \item Tanda tangan digital dari CA.
\end{itemize}

% Rantai kepercayaan sertifikat X.509.
% Dirujuk dari section-03.tex dengan % Rantai kepercayaan sertifikat X.509.
% Dirujuk dari section-03.tex dengan % Rantai kepercayaan sertifikat X.509.
% Dirujuk dari section-03.tex dengan \input{figures/rantai-sertifikat}.
\begin{figure}[htbp]
    \centering
    \begin{tikzpicture}[
        kotak/.style={draw, rounded corners, align=center, font=\small, inner sep=5pt,
                      minimum width=6.4cm, minimum height=1.0cm},
        panah/.style={-{Latex[length=2.4mm]}, thick},
        node distance=1.0cm,
    ]
        \node[kotak, fill=gray!12] (root) {\textbf{Root CA}\\ \scriptsize kunci publik tertanam di peramban dan sistem operasi};
        \node[kotak, fill=blue!6, below=of root] (inter) {\textbf{Intermediate CA}\\ \scriptsize sertifikatnya ditandatangani Root CA};
        \node[kotak, fill=green!10, below=of inter] (server) {\textbf{Sertifikat Server}\\ \scriptsize ditandatangani Intermediate CA};

        \draw[panah] (root) -- node[right, font=\scriptsize, align=left] {menandatangani} (inter);
        \draw[panah] (inter) -- node[right, font=\scriptsize, align=left] {menandatangani} (server);
    \end{tikzpicture}
    \caption{Rantai kepercayaan sertifikat X.509: kepercayaan pada Root CA diteruskan ke sertifikat server melalui Intermediate CA}
    \label{fig:rantai-sertifikat}
\end{figure}
.
\begin{figure}[htbp]
    \centering
    \begin{tikzpicture}[
        kotak/.style={draw, rounded corners, align=center, font=\small, inner sep=5pt,
                      minimum width=6.4cm, minimum height=1.0cm},
        panah/.style={-{Latex[length=2.4mm]}, thick},
        node distance=1.0cm,
    ]
        \node[kotak, fill=gray!12] (root) {\textbf{Root CA}\\ \scriptsize kunci publik tertanam di peramban dan sistem operasi};
        \node[kotak, fill=blue!6, below=of root] (inter) {\textbf{Intermediate CA}\\ \scriptsize sertifikatnya ditandatangani Root CA};
        \node[kotak, fill=green!10, below=of inter] (server) {\textbf{Sertifikat Server}\\ \scriptsize ditandatangani Intermediate CA};

        \draw[panah] (root) -- node[right, font=\scriptsize, align=left] {menandatangani} (inter);
        \draw[panah] (inter) -- node[right, font=\scriptsize, align=left] {menandatangani} (server);
    \end{tikzpicture}
    \caption{Rantai kepercayaan sertifikat X.509: kepercayaan pada Root CA diteruskan ke sertifikat server melalui Intermediate CA}
    \label{fig:rantai-sertifikat}
\end{figure}
.
\begin{figure}[htbp]
    \centering
    \begin{tikzpicture}[
        kotak/.style={draw, rounded corners, align=center, font=\small, inner sep=5pt,
                      minimum width=6.4cm, minimum height=1.0cm},
        panah/.style={-{Latex[length=2.4mm]}, thick},
        node distance=1.0cm,
    ]
        \node[kotak, fill=gray!12] (root) {\textbf{Root CA}\\ \scriptsize kunci publik tertanam di peramban dan sistem operasi};
        \node[kotak, fill=blue!6, below=of root] (inter) {\textbf{Intermediate CA}\\ \scriptsize sertifikatnya ditandatangani Root CA};
        \node[kotak, fill=green!10, below=of inter] (server) {\textbf{Sertifikat Server}\\ \scriptsize ditandatangani Intermediate CA};

        \draw[panah] (root) -- node[right, font=\scriptsize, align=left] {menandatangani} (inter);
        \draw[panah] (inter) -- node[right, font=\scriptsize, align=left] {menandatangani} (server);
    \end{tikzpicture}
    \caption{Rantai kepercayaan sertifikat X.509: kepercayaan pada Root CA diteruskan ke sertifikat server melalui Intermediate CA}
    \label{fig:rantai-sertifikat}
\end{figure}


\subsection{Isi dan Tanda Tangan Sertifikat X.509}
Selain butir-butir di atas, sertifikat X.509 versi 3 memuat \textit{issuer unique ID}, \textit{subject unique ID}, dan blok \textbf{ekstensi} (misalnya \textit{Basic Constraints}, \textit{Key Usage}, \textit{Subject Alternative Name}, dan \textit{CRL Distribution Point}) \citep{stallings2017}. Seluruh data sertifikat ditandatangani CA:
\[ \mathrm{Sig} = \mathrm{Sign}_{K^{\text{CA}}_{\text{priv}}}\bigl(H(\text{data sertifikat})\bigr), \]
sehingga perubahan satu bit pun menggugurkan verifikasi. Nomor seri dan masa berlaku dipakai untuk mencabut atau membatasi keabsahan sertifikat.

\subsection{Rantai Kepercayaan (\textit{Chain of Trust})}
Verifikasi tidak berhenti pada satu sertifikat. Peramban mengikuti rantai: \textbf{Root CA} (sertifikat \textit{self-signed} yang tertanam di \textit{trust store}) $\to$ \textbf{Intermediate CA} (ditandatangani Root) $\to$ \textbf{sertifikat server} (ditandatangani Intermediate). Tiap tautan diverifikasi memakai kunci publik penerbitnya; rantai dinyatakan gagal bila salah satu sertifikat kedaluwarsa, dicabut (melalui CRL/OCSP), atau namanya tidak cocok dengan domain yang diakses \citep{menezes1996}.

